% GENERATED FILE. Edit canonical lessons, then run npm run generate:books.
% canonical-source-hash: fnv1a64:b3cf9740b52ddf9a
% canonical-lessons: SA-C305-uddharati, SA-C305-pranah, SA-C305-sahasam, SA-C305-avatarayati, SA-C305-trayate

\chapter{The Rescue}
\label{ch:sa-the-rescue}
\hlchaptermodality{305}

\begin{chapteropening}
\textbf{By the end of this chapter:} \emph{I can tell how people were rescued.}

Complete the last lesson of chapter 305: I can tell how people were rescued.
\end{chapteropening}

\section[\texorpdfstring{\sk{उद्धरति} (uddharati)}{उद्धरति (uddharati)}]{\sk{उद्धरति} (uddharati) — rescues; pulls clear}

\label{lesson:SA-C305-uddharati}



\begin{quote}
Before the new one: say the Sanskrit for laments, then the Sanskrit for brings out.
\end{quote}

\subsection*{You'll want to know: \sk{उद्धरति}}
\textbf{\sk{उद्धरति}} — \emph{uddharati} — \textquotedblleft{}rescues; pulls clear\textquotedblright{}.

\begin{quote}
\sk{सैनिकाः} \sk{जलात्} \sk{जनान्} \sk{उद्धरन्ति।}
\end{quote}

\begin{quote}
\emph{sainikāḥ jalāt janān uddharanti.}
\end{quote}

\begin{quote}
Soldiers rescue people from the water.
\end{quote}

Past (he …): \sk{सः} \sk{बालं} \sk{जलात्} \sk{उद्धृतवान्}, \emph{saḥ bālaṃ jalāt uddhṛtavān}.

\begin{cousinweb}
\textbf{\sk{उद्धरति}} (\emph{uddharati}) is built from \emph{ud}, \textquotedblleft{}up, out\textquotedblright{}, and the root √hṛ, \textquotedblleft{}to take\textquotedblright{}; MW gives \emph{ud-dhṛ} as to take out, pull out, extricate (MW).
\end{cousinweb}

\begin{grammarlens}[title={What you've built}]
One more word for telling what happened.
\end{grammarlens}

\subsection*{Your turn}
\begin{itemize}
\raggedright
  \item \emph{Say it:} \emph{uddharati}
  \item \emph{Say it:} \emph{uddharati}, once more
  \item \emph{Say it:} \emph{sainikāḥ jalāt janān uddharanti.} — Soldiers rescue people from the water.
  \item \emph{Recall:} say the Sanskrit for laments, then the Sanskrit for brings out, then say \emph{uddharati} again
\end{itemize}

\begin{tcolorbox}[breakable,colback=teal!4,colframe=teal!35!black,title={Before you move on}]
What does \sk{उद्धरति} mean? (\textquotedblleft{}Rescues; pulls clear\textquotedblright{}.) Say it once more.
\end{tcolorbox}
\section[\texorpdfstring{\sk{प्राणाः} (prāṇāḥ)}{प्राणाः (prāṇāḥ)}]{\sk{प्राणाः} (prāṇāḥ) — life; the breath of life}

\label{lesson:SA-C305-pranah}



\begin{quote}
Before the new one: say the Sanskrit for brings out, then the Sanskrit for rescues.
\end{quote}

\subsection*{You'll want to know: \sk{प्राणाः}}
\textbf{\sk{प्राणाः}} — \emph{prāṇāḥ} — \textquotedblleft{}life; the breath of life\textquotedblright{}.

\begin{quote}
\sk{सैनिकः} \sk{बालस्य} \sk{प्राणान्} \sk{रक्षति।}
\end{quote}

\begin{quote}
\emph{sainikaḥ bālasya prāṇān rakṣati.}
\end{quote}

\begin{quote}
The soldier saves the child's life.
\end{quote}

A masculine noun, used in the plural.

\begin{cousinweb}
\textbf{\sk{प्राणाः}} (\emph{prāṇāḥ}) is built from \emph{pra} and the root √an, \textquotedblleft{}to breathe\textquotedblright{}; MW gives \emph{prāṇa} as the breath of life, breath, vitality, and the plural stands for life itself (MW).
\end{cousinweb}

\begin{grammarlens}[title={What you've built}]
One more word for telling what happened.
\end{grammarlens}

\subsection*{Your turn}
\begin{itemize}
\raggedright
  \item \emph{Say it:} \emph{prāṇāḥ}
  \item \emph{Say it:} \emph{prāṇāḥ}, once more
  \item \emph{Say it:} \emph{sainikaḥ bālasya prāṇān rakṣati.} — The soldier saves the child's life.
  \item \emph{Recall:} say the Sanskrit for brings out, then the Sanskrit for rescues, then say \emph{prāṇāḥ} again
\end{itemize}

\begin{tcolorbox}[breakable,colback=teal!4,colframe=teal!35!black,title={Before you move on}]
What does \sk{प्राणाः} mean? (\textquotedblleft{}Life; the breath of life\textquotedblright{}.) Say it once more.
\end{tcolorbox}
\section[\texorpdfstring{\sk{साहसम्} (sāhasam)}{साहसम् (sāhasam)}]{\sk{साहसम्} (sāhasam) — daring; boldness}

\label{lesson:SA-C305-sahasam}



\begin{quote}
Before the new one: say the Sanskrit for rescues, then the Sanskrit for life.
\end{quote}

\subsection*{You'll want to know: \sk{साहसम्}}
\textbf{\sk{साहसम्}} — \emph{sāhasam} — \textquotedblleft{}daring; boldness\textquotedblright{}.

\begin{quote}
\sk{तस्य} \sk{साहसं} \sk{सर्वे} \sk{प्रशंसन्ति।}
\end{quote}

\begin{quote}
\emph{tasya sāhasaṃ sarve praśaṃsanti.}
\end{quote}

\begin{quote}
Everyone praises his daring.
\end{quote}

A neuter noun.

\begin{cousinweb}
\textbf{\sk{साहसम्}} (\emph{sāhasam}) is built on \emph{sahas}, \textquotedblleft{}might, force\textquotedblright{}; MW gives \emph{sāhasa} as boldness, daring, rashness, and also as violence (MW).
\end{cousinweb}

\begin{grammarlens}[title={What you've built}]
One more word for telling what happened.
\end{grammarlens}

\subsection*{Your turn}
\begin{itemize}
\raggedright
  \item \emph{Say it:} \emph{sāhasam}
  \item \emph{Say it:} \emph{sāhasam}, once more
  \item \emph{Say it:} \emph{tasya sāhasaṃ sarve praśaṃsanti.} — Everyone praises his daring.
  \item \emph{Recall:} say the Sanskrit for rescues, then the Sanskrit for life, then say \emph{sāhasam} again
\end{itemize}

\begin{tcolorbox}[breakable,colback=teal!4,colframe=teal!35!black,title={Before you move on}]
What does \sk{साहसम्} mean? (\textquotedblleft{}Daring; boldness\textquotedblright{}.) Say it once more.
\end{tcolorbox}
\section[\texorpdfstring{\sk{अवतारयति} (avatārayati)}{अवतारयति (avatārayati)}]{\sk{अवतारयति} (avatārayati) — brings down; lowers}

\label{lesson:SA-C305-avatarayati}



\begin{quote}
Before the new one: say the Sanskrit for life, then the Sanskrit for daring.
\end{quote}

\subsection*{You'll want to know: \sk{अवतारयति}}
\textbf{\sk{अवतारयति}} — \emph{avatārayati} — \textquotedblleft{}brings down; lowers\textquotedblright{}.

\begin{quote}
\sk{सैनिकाः} \sk{आहतं} \sk{पर्वतात्} \sk{अवतारयन्ति।}
\end{quote}

\begin{quote}
\emph{sainikāḥ āhataṃ parvatāt avatārayanti.}
\end{quote}

\begin{quote}
Soldiers bring the injured man down from the hill.
\end{quote}

Past (he …): \sk{सः} \sk{बालं} \sk{वृक्षात्} \sk{अवतारितवान्}, \emph{saḥ bālaṃ vṛkṣāt avatāritavān}.

\begin{cousinweb}
\textbf{\sk{अवतारयति}} (\emph{avatārayati}) is built from \emph{ava}, \textquotedblleft{}down\textquotedblright{}, and the causative of the root √tṝ, \textquotedblleft{}to cross\textquotedblright{}; MW gives \emph{ava-tārita} as caused to descend, fetched down (MW).
\end{cousinweb}

\begin{grammarlens}[title={What you've built}]
One more word for telling what happened.
\end{grammarlens}

\subsection*{Your turn}
\begin{itemize}
\raggedright
  \item \emph{Say it:} \emph{avatārayati}
  \item \emph{Say it:} \emph{avatārayati}, once more
  \item \emph{Say it:} \emph{sainikāḥ āhataṃ parvatāt avatārayanti.} — Soldiers bring the injured man down from the hill.
  \item \emph{Recall:} say the Sanskrit for life, then the Sanskrit for daring, then say \emph{avatārayati} again
\end{itemize}

\begin{tcolorbox}[breakable,colback=teal!4,colframe=teal!35!black,title={Before you move on}]
What does \sk{अवतारयति} mean? (\textquotedblleft{}Brings down; lowers\textquotedblright{}.) Say it once more.
\end{tcolorbox}
\section[\texorpdfstring{\sk{त्रायते} (trāyate)}{त्रायते (trāyate)}]{\sk{त्रायते} (trāyate) — saves; protects}

\label{lesson:SA-C305-trayate}



\begin{quote}
Before the new one: say the Sanskrit for daring, then the Sanskrit for brings down.
\end{quote}

\subsection*{You'll want to know: \sk{त्रायते}}
\textbf{\sk{त्रायते}} — \emph{trāyate} — \textquotedblleft{}saves; protects\textquotedblright{}.

\begin{quote}
\sk{वैद्यः} \sk{आहतं} \sk{त्रायते।}
\end{quote}

\begin{quote}
\emph{vaidyaḥ āhataṃ trāyate.}
\end{quote}

\begin{quote}
The doctor saves the injured man.
\end{quote}

Past (he …): \sk{सः} \sk{आहतं} \sk{त्रातवान्}, \emph{saḥ āhataṃ trātavān}.

\begin{cousinweb}
\textbf{\sk{त्रायते}} (\emph{trāyate}) comes from the root √trai, \textquotedblleft{}to protect, preserve, rescue\textquotedblright{} (MW).
\end{cousinweb}

\begin{grammarlens}[title={What you've built}]
One more word for telling what happened.
\end{grammarlens}

\subsection*{Your turn}
\begin{itemize}
\raggedright
  \item \emph{Say it:} \emph{trāyate}
  \item \emph{Say it:} \emph{trāyate}, once more
  \item \emph{Say it:} \emph{vaidyaḥ āhataṃ trāyate.} — The doctor saves the injured man.
  \item \emph{Recall:} say the Sanskrit for daring, then the Sanskrit for brings down, then say \emph{trāyate} again
\end{itemize}

\begin{tcolorbox}[breakable,colback=teal!4,colframe=teal!35!black,title={Before you move on}]
What does \sk{त्रायते} mean? (\textquotedblleft{}Saves; protects\textquotedblright{}.) Say it once more.
\end{tcolorbox}
